%% test-japanese.tex
%% =================
%% Typesetting test for the Japanese classes: table notes and page
%% footnotes of tables in Japanese.
%%
%%   lualatex test-japanese.tex   (twice)
%%   platex  -jobname=test-japanese "\def\CLS{jbook}\input test-japanese"
%%   uplatex -jobname=test-japanese "\def\CLS{ujbook}\input test-japanese"
%%   (twice, then dvipdfmx test-japanese)
%% The class is ltjbook by default; with pTeX the class option dvipdfmx
%% is given.
%%
%% What to check:
%%   * TN-RESULT OK, TN-COUNTER OK, and no error or warning.
%%   * In the PDF, the notes of tables 1 and 3 are below the tables, in
%%     their widths, and the footnote of table 2 is at the bottom of
%%     the page.
\ifdefined\CLS\else\def\CLS{ltjbook}\fi
\ifdefined\ptexversion\def\CLSOPT{dvipdfmx}\else\def\CLSOPT{}\fi
\expandafter\documentclass\expandafter[\CLSOPT,a5paper]{\CLS}
\usepackage{longtable,float}
\usepackage{tablenotes}
%% test-common.tex
%% ===============
%% Helpers shared by the test files; \input this file after
%% \usepackage{tblnotes}.
%%
%% \TNfootnote[<kind>]{<text>}
%%   \footnote{<text>} with a check of where its text is typeset.
%%   <kind> (default: table) is one of
%%     table        in the table notes, on the page of its mark;
%%     table-later  in the table notes, on a later page than its mark
%%                  (a longtable broken across pages);
%%     page         a page footnote on the page of its mark;
%%     page-later   a page footnote on a later page than its mark;
%%     page-earlier a page footnote on an earlier page than its mark
%%                  (mode=page: a float that moved to a later page);
%%     box          a footnote of a minipage or a box, on the page of
%%                  its mark.
%% \TNmark{<id>} and \TNtext{<id>}{<kind>}
%%   The two halves of \TNfootnote, for \footnote[<num>],
%%   \footnotemark and \footnotetext: put \TNmark right after the mark
%%   and \TNtext at the beginning of the text.  <id> is any name that
%%   is not a number (numbers are used by \TNfootnote).
%% \TNwidth{<dimen>}
%%   The table notes of the following marks are set in this width
%%   (\TNwidth{} to stop checking the width).
%% \TNcounter{<counter>}{<value>}
%%   Check the value of <counter> at this point.
%%
%% For each mark and each text, the page, the mark (\@thefnmark) and,
%% for a text, where it is set are written to the .aux file at
%% shipout: T (in the table notes of tblnotes), P (in the full
%% \columnwidth of the page, outside the table notes) or B (elsewhere:
%% a minipage or a box).  When the .aux file is read back at
%% \end{document}, every mark is checked:
%%   * its text is typeset exactly once, with the same mark;
%%   * the text is where <kind> says, in the width given by \TNwidth;
%%   * with hyperref, the anchor at the text is that of the mark.
%% Logs "TN-NOTE NG: ..." for each failure, then "TN-RESULT OK (<n>
%% notes)" or "TN-RESULT NG".  Two runs are needed for longtable
%% (column widths).
\makeatletter
\ExplSyntaxOn
\int_new:N  \g__tn_id_int
\seq_new:N  \g__tn_ids_seq
\bool_new:N \g__tn_ng_bool
\dim_new:N  \g__tn_colwd_dim
\tl_new:N   \g__tn_width_tl
\tl_new:N   \l__tn_tmp_tl
\AddToHook { begindocument/end }
  { \dim_gset:Nn \g__tn_colwd_dim { \columnwidth } }

\cs_new_protected:Npn \__tn_ng:n #1
  {
    \bool_gset_true:N \g__tn_ng_bool
    \iow_term:e { TN-NOTE~NG:~#1 }
  }

%% Write to the .aux file at shipout; #1 is frozen now except \thepage
\cs_new_protected:Npn \__tn_shipout:n #1
  {
    \tl_set:Ne \l__tn_tmp_tl {#1}
    \exp_args:NNV \iow_shipout_e:Nn \@auxout \l__tn_tmp_tl
  }
\cs_new:Npn \__tn_frozen_mark:
  { \exp_not:N \exp_not:n { \exp_not:o { \@thefnmark } } }
%% hyperref: the anchor of the mark, and the anchor at the text
%% (hyperref sets \@currentHref to the anchor it puts there)
\cs_new:Npn \__tn_frozen_href:N #1
  {
    \cs_if_exist:NT \Hy@footnote@currentHref
      { \exp_not:N \exp_not:n { \exp_not:o {#1} } }
  }
\cs_new:Npn \__tn_where:
  {
    \bool_if:NTF \l__tblnotes_block_bool
      { T }
      {
        \dim_compare:nNnTF { \hsize } = { \g__tn_colwd_dim }
          { P } { B }
      }
  }

\NewDocumentCommand \TNwidth { m }
  {
    \tl_if_blank:nTF {#1}
      { \tl_gclear:N \g__tn_width_tl }
      { \tl_gset:Ne \g__tn_width_tl { \dim_eval:n {#1} } }
  }
\NewDocumentCommand \TNmark { m }
  {
    \seq_gput_right:Nn \g__tn_ids_seq {#1}
    \__tn_shipout:n
      {
        \exp_not:N \exp_not:N \exp_not:N \TN@m
        {#1} { \exp_not:N \thepage } { \__tn_frozen_mark: }
        { \g__tn_width_tl }
        { \__tn_frozen_href:N \Hy@footnote@currentHref }
      }
  }
\NewDocumentCommand \TNtext { m m }
  {
    \__tn_shipout:n
      {
        \exp_not:N \exp_not:N \exp_not:N \TN@t
        {#1} {#2} { \exp_not:N \thepage } { \__tn_frozen_mark: }
        { \__tn_where: } { \dim_use:N \hsize }
        { \__tn_frozen_href:N \@currentHref }
      }
  }
\NewDocumentCommand \TNfootnote { O{table} +m }
  {
    \int_gincr:N \g__tn_id_int
    \exp_args:NV \__tn_footnote:nnn \g__tn_id_int {#1} {#2}
  }
\cs_new_protected:Npn \__tn_footnote:nnn #1#2#3
  { \footnote { \TNtext {#1} {#2} #3 } \TNmark {#1} }

\NewDocumentCommand \TNcounter { m m }
  {
    \int_compare:nNnTF { \value {#1} } = {#2}
      { \iow_term:e { TN-COUNTER~OK:~#1=#2 } }
      { \__tn_ng:n { counter~#1=\int_use:c { c@#1 },~want~#2 } }
  }

%% The records in the .aux file are ignored when it is read at
%% \begin{document} (they come from the previous run)
\cs_new_eq:NN \TN@m \use_none:nnnnn
\cs_new_eq:NN \TN@t \use_none:nnnnnnn
\cs_new_protected:Npn \__tn_record_mark:nnnnn #1#2#3#4#5
  {
    \tl_gset:cn { g__tn_mhref_#1_tl } {#5}
    \tl_gset:cn { g__tn_mpage_#1_tl } {#2}
    \tl_gset:cn { g__tn_mmark_#1_tl } {#3}
    \tl_gset:cn { g__tn_mwidth_#1_tl } {#4}
  }
\cs_new_protected:Npn \__tn_record_text:nnnnnnn #1#2#3#4#5#6#7
  {
    \tl_gset:cn { g__tn_thref_#1_tl } {#7}
    \int_if_exist:cF { g__tn_tcount_#1_int }
      { \int_new:c { g__tn_tcount_#1_int } }
    \int_gincr:c { g__tn_tcount_#1_int }
    \tl_gset:cn { g__tn_tkind_#1_tl } {#2}
    \tl_gset:cn { g__tn_tpage_#1_tl } {#3}
    \tl_gset:cn { g__tn_tmark_#1_tl } {#4}
    \tl_gset:cn { g__tn_twhere_#1_tl } {#5}
    \tl_gset:cn { g__tn_twidth_#1_tl } {#6}
  }
\AddToHook { enddocument/afterlastpage }
  {
    \cs_gset_eq:NN \TN@m \__tn_record_mark:nnnnn
    \cs_gset_eq:NN \TN@t \__tn_record_text:nnnnnnn
  }
\AddToHook { enddocument/afteraux } { \__tn_check: }

\cs_new_protected:Npn \__tn_check:
  {
    \seq_map_function:NN \g__tn_ids_seq \__tn_check_one:n
    \seq_if_empty:NT \g__tn_ids_seq { \__tn_ng:n { no~footnotes } }
    \bool_if:NTF \g__tn_ng_bool
      { \iow_term:n { TN-RESULT~NG } }
      {
        \iow_term:e
          { TN-RESULT~OK~(\seq_count:N \g__tn_ids_seq\c_space_tl notes) }
      }
  }
\cs_new_protected:Npn \__tn_check_one:n #1
  {
    \tl_if_exist:cTF { g__tn_mpage_#1_tl }
      {
        \int_if_exist:cTF { g__tn_tcount_#1_int }
          {
            \int_compare:nNnTF { \int_use:c { g__tn_tcount_#1_int } } = { 1 }
              { \__tn_check_text:n {#1} }
              {
                \__tn_ng:n
                  { #1:~text~typeset~\int_use:c { g__tn_tcount_#1_int }~times }
              }
          }
          { \__tn_ng:n { #1:~text~lost } }
      }
      { \__tn_ng:n { #1:~mark~not~shipped~out } }
  }
\cs_new_protected:Npn \__tn_check_text:n #1
  {
    \tl_set:Ne \l__tn_tmp_tl
      {
        #1:~\tl_use:c { g__tn_tkind_#1_tl },~
        mark~\tl_to_str:v { g__tn_mmark_#1_tl }~
        on~p.\tl_use:c { g__tn_mpage_#1_tl },~
        text~\tl_to_str:v { g__tn_tmark_#1_tl }~
        on~p.\tl_use:c { g__tn_tpage_#1_tl }~
        (\tl_use:c { g__tn_twhere_#1_tl }~
         \tl_use:c { g__tn_twidth_#1_tl }
         \tl_if_empty:cF { g__tn_mwidth_#1_tl }
           { ,~want~\tl_use:c { g__tn_mwidth_#1_tl } })
        \tl_if_empty:cF { g__tn_mhref_#1_tl }
          {
            ,~anchor~\tl_use:c { g__tn_mhref_#1_tl }~
            ->~\tl_use:c { g__tn_thref_#1_tl }
          }
      }
    \bool_if:nTF
      {
        \tl_if_eq_p:cc { g__tn_mmark_#1_tl } { g__tn_tmark_#1_tl }
        && \__tn_place_p:n {#1}
        && \__tn_width_p:n {#1}
        && \tl_if_eq_p:cc { g__tn_mhref_#1_tl } { g__tn_thref_#1_tl }
      }
      { \iow_term:e { TN-NOTE~OK:~\l__tn_tmp_tl } }
      { \__tn_ng:n { \l__tn_tmp_tl } }
  }
%% #1: id, #2: the code of the place, #3: the relation of the pages
\cs_new:Npn \__tn_is:nnN #1#2#3
  {
    \str_if_eq_p:ee { \tl_use:c { g__tn_twhere_#1_tl } } {#2}
    && \int_compare_p:nNn { \tl_use:c { g__tn_tpage_#1_tl } }
         #3 { \tl_use:c { g__tn_mpage_#1_tl } }
  }
\prg_new_conditional:Npnn \__tn_place:n #1 { p }
  {
    \str_case_e:nnF { \tl_use:c { g__tn_tkind_#1_tl } }
      {
        { table }
          { \bool_if:nTF { \__tn_is:nnN {#1} { T } = } { \prg_return_true: } { \prg_return_false: } }
        { table-later }
          { \bool_if:nTF { \__tn_is:nnN {#1} { T } > } { \prg_return_true: } { \prg_return_false: } }
        { page }
          { \bool_if:nTF { \__tn_is:nnN {#1} { P } = } { \prg_return_true: } { \prg_return_false: } }
        { page-later }
          { \bool_if:nTF { \__tn_is:nnN {#1} { P } > } { \prg_return_true: } { \prg_return_false: } }
        { page-earlier }
          { \bool_if:nTF { \__tn_is:nnN {#1} { P } < } { \prg_return_true: } { \prg_return_false: } }
        { box }
          { \bool_if:nTF { \__tn_is:nnN {#1} { B } = } { \prg_return_true: } { \prg_return_false: } }
      }
      { \prg_return_false: }
  }
\prg_new_conditional:Npnn \__tn_width:n #1 { p }
  {
    \tl_if_empty:cTF { g__tn_mwidth_#1_tl }
      { \prg_return_true: }
      {
        \dim_compare:nNnTF
          { \dim_abs:n { \tl_use:c { g__tn_twidth_#1_tl }
                         - \tl_use:c { g__tn_mwidth_#1_tl } } }
          < { 0.01pt }
          { \prg_return_true: } { \prg_return_false: }
      }
  }
\ExplSyntaxOff

%% \TNfiller{<n>}: a paragraph of <n> filler sentences
\newcommand\TN@filler{%
  This is a filler sentence of plain English text for the test. }
\newcommand\TNfiller[1]{%
  \count@=0
  \loop\advance\count@ by 1 \TN@filler\ifnum\count@<#1 \repeat\par}
%% \TNrows{<from>}{<to>}: rows "R<n> & x \\" of a two-column table
\newcount\TN@row
\newcommand\TNrows[2]{%
  \global\TN@row=#1\relax
  \TN@rows{#2}}
\newcommand\TN@rows[1]{%
  \ifnum\TN@row>#1\relax
    \expandafter\@gobble
  \else
    \expandafter\@firstofone
  \fi
  {R\the\TN@row & x \\ \noalign{\global\advance\TN@row\@ne}\TN@rows{#1}}}
\makeatother
\endinput

\begin{document}

\section*{表中の脚注のテスト}

これは通常の本文です。本文の脚注\TNfootnote[page]{本文中の通常の脚注です。ページ下部に出力されます。}も書けます。

\begin{table}[ht]
\centering
\caption{表注（既定）}
\TNwidth{4cm+5cm+4\tabcolsep}
\begin{tabular}{|p{4cm}|p{5cm}|}
\hline
項目\TNfootnote{表注は表の下に、表の幅で出力されます。} &
説明\TNfootnote{番号はmpfootnoteカウンタで、表ごとに振り直されます。} \\
\hline
\end{tabular}
\TNwidth{}
\end{table}

\begin{table}[H]
\centering
\tablenotessetup{mode=page}
\caption{本文の脚注として出力}
\begin{tabular}{|l|l|}
\hline
項目\TNfootnote[page]{mode=pageでは、本文の脚注と通し番号でページ下部に出力されます。} & 説明 \\
\hline
\end{tabular}
\end{table}

\begin{longtable}{|l|p{6cm}|}
\caption{longtableの表注}\\
\hline
項目\TNfootnote{longtableの表注も、表の下に出力されます。} &
長い表\TNfootnote{p列の脚注も回収されます。} \\
\hline
\end{longtable}

最後に本文の脚注\TNfootnote[page]{最後の脚注です。}を置きます。

\TNcounter{footnote}{3}
\end{document}
